\documentclass[12pt,a4paper]{article}
\usepackage[spanish,es-nodecimaldot]{babel}
\usepackage{fontspec}
% Instalar Times New Roman y compilar con XeLaTeX para usar la fuente solicitada.
\setmainfont{Times New Roman}
\usepackage[margin=2.3cm]{geometry}
\usepackage{longtable,array,booktabs,fvextra,enumitem,fancyhdr,needspace}
\usepackage[hidelinks,unicode]{hyperref}
\usepackage{xurl}
\setlength{\parindent}{0pt}
\setlength{\parskip}{6pt}
\setlength{\emergencystretch}{3em}
\setlist{nosep,leftmargin=1.5em,topsep=5pt}
\pagestyle{fancy}\fancyhf{}
\fancyhead[L]{\small IDENTITY ACCESS LAB}\fancyhead[R]{\small Manual de laboratorio}
\fancyfoot[C]{\thepage}
\setlength{\headheight}{15pt}
\DefineVerbatimEnvironment{Comandos}{Verbatim}{breaklines,breakanywhere,fontsize=\small,fontfamily=\rmdefault,frame=single,framesep=5pt}
\begin{document}
{\LARGE\bfseries IAM: OneLogin, permisos y ciclo de vida\par}\vspace{6pt}
2 de octubre de 2026\par\hrule\vspace{10pt}
\Needspace{5\baselineskip}\section*{Objetivo y reparto}
Conectar OneLogin con la aplicación para demostrar identidad real, MFA, permisos por puesto y altas/bajas automáticas. No sustituirlo por otro proveedor ni por una pantalla de login creada por el grupo.\par
IAM tiene cinco puestos: coordinación de roles, conexión OneLogin, MFA y ciclo de vida, pruebas de aprovisionamiento y documentación técnica. Asignar cada función en el control privado.\par
Laptop IAM 1: administración. IAM 2: navegador de pruebas, con perfiles separados para ingeniero, auditor y usuario sin permiso.\par
\Needspace{5\baselineskip}\section*{Antes de iniciar}
Necesitan tenant OneLogin, cuenta administrativa autorizada, aplicación HTTPS activa y capacidad de aprovisionamiento en el plan disponible. Confirmar estas funciones al principio. El acceso de prueba no garantiza todas las funciones: si faltan, solicitar a la institución habilitación o el entorno adecuado.\par
Recibir de TI APP\_BASE\_URL y SCIM\_TOKEN. Estos son datos de configuración; el token no debe aparecer en capturas. Las operaciones se hacen sobre cuentas de laboratorio.\par
\Needspace{5\baselineskip}\section*{Paso 1. Crear usuarios de prueba}
\begin{enumerate}
\item Abrir el portal administrativo del tenant.
\item Ir a Users \(\rightarrow\) Users y crear las identidades del recorrido según el registro de RRHH.
\item Usar el mismo correo en OneLogin, Microsoft 365 y el registro. No es sincronización automática: son cuentas relacionadas por acuerdo.
\item Registrar LAB-INGENIERO para ingeniero, LAB-ADMIN para administrador y LAB-AUDITOR para auditor.
\item Preparar también una identidad sin permiso y otra temporal para alta/baja, si los recursos lo permiten.
\item Hacer que cada usuario establezca su contraseña de laboratorio y registrar quién es responsable de su método MFA.
\end{enumerate}
\Needspace{5\baselineskip}\section*{Paso 2. Definir atributos y roles}
Crear campos de usuario \textnormal{lab\_\allowbreak{}role} y \textnormal{emp\_\allowbreak{}id} mediante la opción de campos personalizados del tenant. Registrar sus nombres cortos exactamente así.\par
{\small\renewcommand{\arraystretch}{1.2}
\begin{longtable}{p{5.027cm}p{5.027cm}p{5.027cm}}
\toprule
\textbf{Identidad} & \textbf{lab\_role} & \textbf{Uso} \\ \midrule\endfirsthead
\textbf{Identidad} & \textbf{lab\_role} & \textbf{Uso} \\ \midrule\endhead
LAB-INGENIERO & mantenimiento & Consultar y ejecutar mantenimiento \\
LAB-ADMIN & admin & Consultar, mantener, administrar y ver registros \\
LAB-AUDITOR & auditor & Consultar y ver registros; sin mantenimiento \\
Usuario de consulta & consulta & Solo consultar \\
Usuario sin acceso & sin\_acceso & Ninguna función \\
\bottomrule\end{longtable}}
El valor \textnormal{title} que recibirá la aplicación por SCIM transporta el código \textnormal{lab\_\allowbreak{}role} en este ejercicio. No es el nombre humano del puesto. El puesto completo permanece en RRHH.\par
Crear roles en Users \(\rightarrow\) Roles: \textnormal{LAB-Mantenimiento}, \textnormal{LAB-Administracion}, \textnormal{LAB-Auditoria} y, si se usa, \textnormal{LAB-Consulta}. Un rol OneLogin agrupa aplicaciones permitidas; la aplicación comprueba las funciones internas usando el atributo sincronizado.\par
\Needspace{5\baselineskip}\section*{Paso 3. Crear conector de inicio de sesión}
\begin{enumerate}
\item Ir a Applications/Apps \(\rightarrow\) Add Apps y localizar el conector OpenID Connect (OIDC) de OneLogin.
\item Nombrarlo \textnormal{IDENTITY ACCESS LAB - Servidor}.
\item Configurar el flujo Authorization Code para una aplicación con servidor. No usar Password Grant ni un secreto dentro del navegador.
\item Registrar callback exacto: \textnormal{APP\_\allowbreak{}BASE\_\allowbreak{}URL/\allowbreak{}auth/\allowbreak{}callback}.
\item Registrar URL de inicio como \textnormal{APP\_\allowbreak{}BASE\_\allowbreak{}URL/\allowbreak{}login} si el conector solicita Login URL.
\item Copiar Client ID y Client Secret al responsable de TI para \textnormal{.env}.
\item Copiar la dirección de configuración del proveedor OIDC. Para API v2 suele usar \textnormal{/\allowbreak{}oidc/\allowbreak{}2/\allowbreak{}.well-known/\allowbreak{}openid-configuration}; confirmar con el tenant.
\item Verificar método de autenticación del cliente. Authlib usa el método compatible configurado; si el tenant exige POST y rechaza Basic, TI debe configurar \textnormal{token\_\allowbreak{}endpoint\_\allowbreak{}auth\_\allowbreak{}method=\allowbreak{}'client\_\allowbreak{}secret\_\allowbreak{}post'} en \textnormal{client\_\allowbreak{}kwargs}.
\item Asociar esta aplicación a los roles de laboratorio autorizados.
\end{enumerate}
Resultado: al pulsar login en la aplicación, aparece la página real de OneLogin. Todavía puede haber rechazo al volver si la cuenta no ha sido creada por SCIM.\par
\Needspace{5\baselineskip}\section*{Paso 4. Configurar MFA OneLogin}
\begin{enumerate}
\item Abrir Security \(\rightarrow\) Authentication Factors y habilitar OneLogin Protect u otro factor aceptado por el tenant.
\item Crear una política de usuario para el laboratorio en Security \(\rightarrow\) Policies.
\item En su apartado MFA, exigir verificación en cada login para la demostración, cuando aparezca la opción \textnormal{At every login}.
\item Asociar la política a los usuarios de prueba mediante el grupo o configuración de usuario correspondiente.
\item En IAM 2, iniciar sesión con el ingeniero y registrar el factor en su teléfono.
\item Cerrar sesión de OneLogin y aplicación, abrir un perfil limpio y repetir.
\item Guardar evidencia del desafío y resultado. No fotografiar QR de inscripción ni códigos secretos.
\end{enumerate}
El método puede estar en un teléfono distinto del usado como cámara. Si comparten dispositivo, comprobar que cambiar a la app MFA no detiene la cámara antes del ensayo.\par
\Needspace{5\baselineskip}\section*{Paso 5. Crear conector de altas y bajas}
El servidor incluido implementa el perfil de práctica SCIM Core 1.0 del conector documentado de OneLogin. No se presenta como una implementación completa de SCIM 2.0.\par
\begin{enumerate}
\item Añadir \textnormal{SCIM Provisioner with SAML (Core Schema)} y nombrarlo \textnormal{IDENTITY ACCESS LAB - Provisioning}.
\item Usar solo su función de aprovisionamiento. El login sigue siendo OIDC; no configurar una URL SAML inexistente para aparentar SSO.
\item En Configuration, colocar SCIM Base URL: \textnormal{APP\_\allowbreak{}BASE\_\allowbreak{}URL/\allowbreak{}scim/\allowbreak{}v1}.
\item Colocar el token que TI guardó como SCIM\_TOKEN.
\item Pegar el archivo \textnormal{Implementacion/\allowbreak{}app/\allowbreak{}scim-template.json} como plantilla.
\item En Parameters, mapear \textnormal{scimusername} al correo, \textnormal{lab\_\allowbreak{}role} al campo personalizado correspondiente y \textnormal{emp\_\allowbreak{}id} al ID de empleado. Si faltan parámetros, agregarlos y habilitar su inclusión en el aprovisionamiento según la interfaz.
\item No activar sincronización de grupos: el ejemplo sincroniza usuarios y su rol, no grupos.
\item Habilitar la conexión y revisar API Status. Si falla, comprobar la URL HTTPS, token y registros del servidor.
\item En Provisioning, habilitar creación y actualización de usuarios. Para la baja, seleccionar suspensión/desactivación de cuenta y probar qué envío realiza el conector.
\item Asociar este conector a los mismos roles de laboratorio que el OIDC.
\end{enumerate}
La aplicación admite GET/POST de usuarios, consulta por \textnormal{userName}, actualización PUT/PATCH y eliminación. Acepta \textnormal{active:false} y niega acciones a cuentas inactivas. Si se activa una función distinta, como grupos, hay que implementar y probar esa extensión antes de declararla soportada.\par
\Needspace{5\baselineskip}\section*{Paso 6. Automatizar asignación}
En Users \(\rightarrow\) Mappings, crear condiciones según el atributo \textnormal{lab\_\allowbreak{}role} y asignar el rol OneLogin correspondiente. Ejemplo: si \textnormal{lab\_\allowbreak{}role} es \textnormal{mantenimiento}, asignar \textnormal{LAB-Mantenimiento}, que incluye ambos conectores.\par
\begin{enumerate}
\item Guardar la regla.
\item Aplicarla al usuario de prueba y comprobar el rol resultante.
\item Revisar opciones de retirar roles que dejan de cumplir condiciones. No asumir que una regla solo de alta retira automáticamente un permiso viejo.
\item Para el primer ensayo, se puede aprobar manualmente una operación pendiente en Users \(\rightarrow\) Provisioning.
\item Para demostrar automatización completa del laboratorio, configurar las aprobaciones permitidas en ese conector de prueba y repetir una nueva alta sin crear la cuenta manualmente en la aplicación.
\item Registrar si la acción fue automática o si requirió aprobación administrativa. La aprobación no equivale a una creación manual de cuenta, pero debe explicarse.
\end{enumerate}
\Needspace{5\baselineskip}\section*{Paso 7. Probar alta automática}
\begin{enumerate}
\item Confirmar que el correo temporal no existe en la aplicación.
\item Crear o activar el usuario en OneLogin con sus atributos.
\item Dejar que Mapping asigne su rol o realizar la asignación de rol acordada.
\item Revisar Users \(\rightarrow\) Provisioning hasta que el evento esté completado.
\item TI confirma que se creó el usuario local y SOC guarda el evento \textnormal{SCIM\_\allowbreak{}CREATE}.
\item Iniciar sesión por OIDC y ejecutar consulta y mantenimiento según rol.
\end{enumerate}
Una cuenta creada a mano en SQLite o por un curl de prueba no demuestra alta automática desde OneLogin.\par
\Needspace{5\baselineskip}\section*{Paso 8. Probar autorización}
{\small\renewcommand{\arraystretch}{1.2}
\begin{longtable}{p{5.027cm}p{5.027cm}p{5.027cm}}
\toprule
\textbf{Usuario} & \textbf{Acción} & \textbf{Resultado esperado} \\ \midrule\endfirsthead
\textbf{Usuario} & \textbf{Acción} & \textbf{Resultado esperado} \\ \midrule\endhead
Ingeniero & Consulta & Permitido \\
Ingeniero & Mantenimiento & Permitido \\
Ingeniero & Administración & 403, denegado \\
Auditor & Consulta y registros & Permitido \\
Auditor & Mantenimiento & 403, denegado \\
Sin aplicación asignada & Login/aplicación & Rechazo según control de OneLogin o aplicación \\
\bottomrule\end{longtable}}
Probar con botón o petición real, no solo mirando la pantalla. Para el auditor, el botón mantenimiento se deja visible intencionalmente y el servidor debe rechazarlo.\par
\Needspace{5\baselineskip}\section*{Paso 9. Probar cambio y baja}
\begin{enumerate}
\item Cambiar \textnormal{lab\_\allowbreak{}role} de mantenimiento a auditor en OneLogin.
\item Comprobar actualización local por SCIM.
\item Sin cerrar la sesión anterior, intentar mantenimiento: debe ser rechazado después de recibir la actualización.
\item Probar consulta y registros permitidos.
\item Desactivar el usuario según la solicitud de RRHH.
\item Esperar y registrar el evento de suspensión recibido, con su hora real.
\item Probar acción en sesión previa y nuevo login. La aplicación consulta el estado local en cada acción; tras recibir la baja, la sesión anterior no permite ejecutar funciones.
\item Verificar baja en Microsoft 365 y RFID con los otros departamentos. Esas acciones no ocurren automáticamente por configurar SCIM de esta aplicación.
\end{enumerate}
\Needspace{5\baselineskip}\section*{Comprobaciones técnicas}
TI puede comprobar \textnormal{/\allowbreak{}health} sin credenciales. Una consulta SCIM sin token debe devolver 401:\par
\Needspace{4\baselineskip}
\begin{Comandos}
curl -i http://127.0.0.1:8000/scim/v1/Users
\end{Comandos}
En la terminal donde TI cargó \textnormal{.env}, consultar un correo de laboratorio:\par
\Needspace{6\baselineskip}
\begin{Comandos}
curl --get http://127.0.0.1:8000/scim/v1/Users \
  -H "Authorization: Bearer ${SCIM_TOKEN}" \
  --data-urlencode 'filter=userName eq "REEMPLAZAR_CORREO"'
\end{Comandos}
No pegar el token literal en capturas. Si devuelve \textnormal{totalResults:0}, la cuenta no llegó o se usó otro correo. Si hay 401, revisar token. Si la conexión da 502, revisar Gunicorn y túnel.\par
\Needspace{5\baselineskip}\section*{Entregan}
Tabla de roles, configuración OIDC, configuración SCIM sin secretos, reglas, evidencia MFA, altas, cambios, bajas y eventos OneLogin. Escribir cuáles acciones fueron automáticas y cuáles manuales.\par
\Needspace{5\baselineskip}\section*{Referencias y videos}
\begin{itemize}
\item \href{https://developers.onelogin.com/docs/openid-connect/connect-to-onelogin/}{OneLogin: conectar aplicación OIDC}.
\item \href{https://developers.onelogin.com/docs/openid-connect/api/provider-config/}{OneLogin: configuración del proveedor}.
\item \href{https://developers.onelogin.com/docs/scim/define-user-schema/}{OneLogin: esquema de usuarios SCIM}.
\item \href{https://developers.onelogin.com/docs/scim/implement-scim-api/}{OneLogin: implementar API SCIM}.
\item \href{https://developers.onelogin.com/docs/scim/create-app/}{OneLogin: crear aplicación SCIM de prueba}.
\item \href{https://www.onelogin.com/getting-started/free-trial-plan/add-mfa}{OneLogin: agregar MFA}.
\item \href{https://www.onelogin.com/video/groups-roles-and-mappings-oh-my}{Video oficial: grupos, roles y mappings}.
\item \href{https://www.onelogin.com/resource-center/videos/getting-started-with-identity-lifecycle-management-pt-2}{Video oficial: ciclo de vida, parte 2}.
\end{itemize}
\end{document}
